\subsection{Read-out adaptation lowers MAE without reliable position tracking}
\label{sec:forecast-result}

The original decoder omits the wall from predicted latents within two steps for all three models,
although it renders the wall from real encodings. To examine this transfer, we fit identical affine
ridge read-outs on real future encodings, forecasts, or copied current encodings and evaluate their
full cross-domain transfer. Training and evaluation use separate archives. Both predeclared
train-only scaling schemes, four horizons, all models, and copied-current and constant baselines
are retained in Appendix~\ref{app:forecast}.

Refitting reduces model 2's eight-step wall-position MAE from $16.47$ to $9.11$ pixels on the
wall-present cohort. However, the forecast-fitted probe has RMSE $18.65$ pixels, worse than $13.49$
for the constant training-mean baseline. Its correlation with true position is only $0.246$.
In the prespecified current-pre-reveal continuing slice ($n=716$), refitting reduces MAE from
$10.02$ to $6.16$ pixels, a paired gain of $3.85$ $[3.48,4.22]$. Yet its within-slice correlation
is negative ($-0.175$). A post-hoc check using only the matching training cohort gives constant
mean and median baselines of $4.77$ and $4.64$ pixels, both better than the adapted probe. The
predeclared global-mean baseline gives $14.24$ pixels on this slice, while the matching-cohort
constants outperform the adapted probe. The adapted probe's advantage over a copied-current probe
also has an interval containing zero. These measurements show changed read-out error, with a coarse
association in one model; they do not demonstrate reliable position recovery. Models 0 and 1 have
weaker real-latent controls and forecast-fitted errors near the constant baseline.

The forecast test therefore leaves the representation question unresolved at the precision needed
for control. A missing decoded wall establishes failure of that decoder on forecasts. Refitting a
state probe, by itself, does not turn a favorable conditional MAE into evidence of accurate dynamics.
